\section{From localization to parity}
\label{sec:parity}

We finish by relating the output probability to transitions between
product states. All localization estimates have already been proved;
no assumption about coherent or clean computation is needed here.
For the broader Fourier-concentration setting, see~\cite{GGJ2026};
the character-orthogonality identity needed here is proved directly.

\begin{proof}[Proof of \Cref{thm:parity}]
Suppose that a circuit as in the theorem exists on \(n\) input qubits
and \(N-n\) initially zero ancillas, with \(n\le N\le n^c\).
Retain any discarded qubits as idle wires and write \(W\) for the
resulting unitary. By \Cref{sec:localization}, \(W\in\cU_d(N)\)
after padding with identity layers.
Let \(O\) project the output qubit onto \(\ket0\), and put
\[
 B=W^\dagger OW,\qquad
 p_0(x)=\bra{x,0}B\ket{x,0}.
\]
The Born rule identifies \(p_0(x)\) with the probability of output
zero, irrespective of the final states of the other qubits.
Here and below the last zero denotes all initially zero ancillas.

For \(h\in\{0,1\}^n\), define
\[
 \ket{\xi_h}
 =\bigotimes_{j=1}^n
   \frac{\ket0+(-1)^{h_j}\ket1}{\sqrt2},
 \qquad
 \bar h=(1-h_1,\ldots,1-h_n).
\]
Let \(M_h\) be the count on the input qubits whose reference vectors
are the factors of \(\xi_h\). The vectors
\(\ket{\xi_h,0}\) and \(\ket{\xi_{\bar h},0}\) have exact counts
zero and \(n\), respectively.
Choose fixed \(0<s<t<1/c\). For all sufficiently large \(n\),
\(N^t\le n^{ct}<n\). Since \(O\) is a product projection and
\(W^\dagger\in\cU_d(N)\), \Cref{thm:localization} gives
\begin{equation}
\label{eq:complementary-entry}
 \left|\bra{\xi_{\bar h},0}B\ket{\xi_h,0}\right|
 \le e^{-N^s}
 \qquad\text{for every }h.
\end{equation}
Its uniformity in the count makes the same lower bound on \(N\)
valid for all \(h\).

Averaging these complementary matrix entries gives the parity Fourier
coefficient of $p_0$, namely $2^{-n}\sum_x(-1)^{|x|}p_0(x)$,
where $|x|=\sum_jx_j$. To see this, note that
\(\langle x\mid\xi_h\rangle=2^{-n/2}(-1)^{h\cdot x}\), where
the dot product is taken modulo two. Also,
\((-1)^{\bar h\cdot x}=(-1)^{|x|+h\cdot x}\).
Character orthogonality therefore gives
\begin{align}
 \E_h\bra{\xi_{\bar h},0}B\ket{\xi_h,0}
 &=2^{-n}\sum_{x,y}(-1)^{|x|}
   \E_h(-1)^{h\cdot(x+y)}\bra{x,0}B\ket{y,0}\notag\\
 &=2^{-n}\sum_x(-1)^{|x|}p_0(x).
\label{eq:parity-average}
\end{align}
The addition in \(x+y\) in the character is modulo two.
For even \(x\), the success guarantee gives
\(p_0(x)\ge1/2+\varepsilon\); for odd \(x\), it gives
\(p_0(x)\le1/2-\varepsilon\).
The even and odd classes have equal size, so the final expression
in \eqref{eq:parity-average} is at least \(\varepsilon\).
But \eqref{eq:complementary-entry} bounds its absolute value by
\(e^{-N^s}\), which tends to zero because \(N\ge n\).
This contradiction proves the theorem.
\end{proof}

The role of a polynomial qubit bound is solely the choice of
\(0<s<t<1/c\). The localization theorem itself is uniform on
\(N\) qubits and does not assume any relation between \(N\) and an
input length. The proof gives no claim of the same lower bound
when the ancillary register is allowed to be superpolynomial.

